\chapter[Analyse statistique et probabiliste des écarts de gravité]{Analyse statistique et probabiliste\\des écarts de gravité}
\label{chap:statistiques}

Le notebook 05 prolonge l'exploration par des écarts absolus, des rapports de proportions
et des intervalles bootstrap. La question n'est pas seulement de savoir si deux
pourcentages diffèrent~: il faut préciser la quantité estimée, son dénominateur et le
processus de répétition auquel son intervalle se rapporte.

\section{Probabilités conditionnelles et variable de Bernoulli}
Pour un usager à gravité connue, $Y$ vaut 1 en cas de décès ou d'hospitalisation et 0
sinon. Pour un usager tiré au hasard dans la population définie, $Y$ est une variable
de Bernoulli de paramètre $p$, avec $E(Y)=p$ et $\operatorname{Var}(Y)=p(1-p)$.
Cette description marginale n'impose pas que les usagers d'un même accident soient
indépendants.

Notons $S$ le périmètre des usagers d'accidents corporels enregistrés en métropole en
2019-2022, avec gravité connue. Pour un groupe $G$, la quantité étudiée est
\begin{equation}
 p_G=P(Y=1\mid G,S),\qquad
 \widehat p_G=\frac{n_{G,1}}{n_{G,0}+n_{G,1}},
 \label{eq:proportion}
\end{equation}
où $n_{G,1}$ et $n_{G,0}$ sont les effectifs positifs et négatifs connus du groupe.
L'estimateur est pondéré par les usagers, pas par un poids égal attribué à chaque accident.

\section{Contrastes retenus et mesure de leur ampleur}
Quatre comparaisons sont fixées à partir de l'exploration~: 75 ans ou plus contre
35-44 ans, piétons contre conducteurs, nuit sans éclairage public contre plein jour,
et hors agglomération contre en agglomération. L'orientation A / R reste identique
dans les tableaux, les graphiques et les réplications.

Deux mesures sont complémentaires~:
\begin{equation}
 \Delta=100(p_A-p_R)\quad\text{en points de pourcentage},\qquad
 R_p=\frac{p_A}{p_R}\quad\text{sans unité}.
 \label{eq:contrastes}
\end{equation}
Un écart entre 32,6\,\% et 16,6\,\% représente environ 16 points de pourcentage et un
rapport de 1,97. Ce rapport est un rapport de proportions conditionnelles, pas un
rapport de risques d'accident par trajet, ni un rapport de cotes (\textit{odds ratio}).
Si la proportion de référence est nulle, le rapport reste indéfini.

Dans chaque contraste, A et R sont disjoints. En revanche, les groupes de dimensions
différentes se recoupent~: un piéton peut aussi avoir plus de 75 ans et être impliqué
hors agglomération. Les comptages des neuf groupes, ensemble inclus, ne doivent donc
pas être additionnés pour reconstruire la population.

\section{Quel sens donner à un intervalle ici~?}
Les proportions sont des descriptions exactes des enregistrements disponibles. BAAC
n'est pas un échantillon aléatoire simple de tous les déplacements. Les intervalles
ne sont donc pas une marge d'erreur garantie sur l'ensemble des accidents réels.

Le cadre retenu est une répétition hypothétique~: si des accidents comparables étaient
enregistrés à nouveau dans les mêmes années et avec des pratiques de collecte
comparables, quelle variabilité nos statistiques pourraient-elles présenter~? Les
accidents sont supposés approximativement indépendants au sein d'une année et leur
distribution empirique représente le processus envisagé. Les années sont fixes et
les personnes restent liées à l'intérieur de chaque accident.

Ce cadre ne modélise pas une dépendance spatiale ou temporelle supplémentaire entre
accidents. Il ne corrige ni la sous-déclaration, ni les erreurs de renseignement, ni
les gravités inconnues. La prudence d'interprétation reste nécessaire même avec un
grand effectif et un intervalle étroit.

\section{Bootstrap par accident, stratifié par année}
Le rééchantillonnage de groupes entiers et la sélection séparée dans les strates sont
des mécanismes décrits par le manuel \textit{bsample}~\cite{stata-bsample}. Notre choix
de l'accident comme groupe et de l'année comme strate répond à l'organisation de BAAC.
La référence décrit le mécanisme, sans valider à elle seule toutes ses hypothèses ici.

Un exemple fictif rend la règle concrète. Trois accidents A, B et C comportent
respectivement 2, 1 et 3 usagers, dont 1, 0 et 2 cas positifs. La proportion observée
est $3/6=50\,\%$. Un tirage A, A, C conserve les deux copies de A~: il comprend sept
usagers et quatre positifs, soit $4/7$. Le nombre de tirages d'accidents est inchangé,
mais le nombre d'usagers varie. Éliminer les tirages répétés ou conserver seulement
un usager par accident changerait la méthode.

Pour chaque réplication réelle, autant d'accidents que ceux observés en 2019 sont
tirés avec remise, puis l'opération est répétée séparément pour 2020, 2021 et 2022.
Les contributions sont regroupées pour calculer les proportions de la période.
Les 207\,057 accidents sont conservés dans la base de tirage, y compris ceux dont
la contribution est nulle pour un groupe donné ou pour les gravités connues.

L'implémentation utilise une matrice de comptages par accident~: pour chacun des neuf
groupes, elle stocke les gravités connues, les cas positifs principaux et les décès.
Les 27 colonnes suffisent aux statistiques calculées. La somme des comptages des
accidents tirés est équivalente à la copie de leurs usagers, avec leurs répétitions.
Cette équivalence a été vérifiée sur un exemple indépendant.

Le nombre de réplications est $B=2\,000$ et la graine est 20261016. Chaque année utilise
\code{numpy.default_rng} avec une initialisation fondée sur la graine et l'année.
Les mêmes tirages servent à tous les groupes et aux deux définitions de gravité~:
les liens entre personnes et entre statistiques sont ainsi conservés. Les intervalles
des écarts sont calculés sur les écarts répliqués, sans soustraire les bornes de deux
intervalles de proportions séparés.

\subsection{Construction des bornes percentiles}
Pour une statistique $T$, les réplications $T_1^*,\ldots,T_B^*$ donnent l'intervalle
percentile nominal à 95\,\%~\cite{nist-bootstrap}~:
\begin{equation}
 \left[q_{0,025}(T^*),\ q_{0,975}(T^*)\right].
 \label{eq:ic}
\end{equation}
Les quantiles utilisent une interpolation linéaire. Il s'agit d'une approximation
simple non corrigée du biais. Le niveau de 95\,\% concerne la procédure dans un cadre
de répétition~; il n'attribue pas une probabilité de 95\,\% à un paramètre fixé à
l'intérieur de l'intervalle déjà calculé.

Les intervalles sont marginaux, un par statistique, et ne garantissent pas une
couverture simultanée de toutes les comparaisons. Si une réplication rend une
statistique indéfinie, le programme signale le nombre de valeurs valides et laisse
l'intervalle manquant, sans écarter silencieusement cette réplication. Tous les
contrastes présentés disposent ici de 2\,000 valeurs valides.

\section{Résultats pour décès ou hospitalisation}
Le tableau~\ref{tab:deces_hospitalisation} et la figure~\ref{fig:forest-principal}
présentent les résultats regroupés. Les quatre écarts sont positifs et leurs
intervalles restent au-dessus de zéro dans le cadre de répétition retenu.
\begin{table}[htbp]
\centering\small
\renewcommand{\arraystretch}{1.18}
\caption{Écarts et rapports : décès ou hospitalisation, 2019-2022.}\label{tab:deces_hospitalisation}
\begin{tabular*}{\textwidth}{@{\extracolsep{\fill}}lrrrr@{}}
\toprule
Comparaison A / R & Écart (points) & IC à 95\% & Rapport & IC à 95\% \\
\midrule
75+ / 35-44 ans & 19,35 & [18,69 ; 20,04] & 2,39 & [2,33 ; 2,46] \\
Piétons / conducteurs & 16,03 & [15,52 ; 16,54] & 1,97 & [1,93 ; 2,00] \\
Nuit sans éclairage / jour & 11,85 & [11,35 ; 12,38] & 1,70 & [1,67 ; 1,73] \\
Hors / en agglomération & 13,02 & [12,76 ; 13,28] & 2,01 & [1,98 ; 2,04] \\
\bottomrule
\end{tabular*}
\sourcecalcul{notebook 05 ; intervalles percentiles marginaux, 2 000 réplications d'accidents stratifiées par année}
\end{table}

\figurerapport{05_01_contrastes_cible_principale.png}{Écarts en points et rapports de proportions pour la cible principale. Les références sont respectivement zéro et un.}{fig:forest-principal}

Pour les piétons, la proportion atteint 32,60\,\%, contre 16,57\,\% pour les conducteurs.
L'écart est de 16,03 points, avec un intervalle [15,52~; 16,54], et le rapport de 1,97,
avec un intervalle [1,93~; 2,00]. Pour l'âge, l'écart est plus grand, à 19,35 points,
mais il concerne deux tranches choisies après l'exploration. Les rapports de 1,70 pour
la lumière et 2,01 pour l'agglomération résument des différences brutes.

Ces résultats ne désignent pas quatre effets indépendants. Une partie de leurs écarts
peut refléter des compositions communes en âge, véhicule ou contexte. Le fait qu'un
intervalle exclue zéro ne transforme pas l'association en effet causal et ne résout
pas le caractère exploratoire de la sélection des comparaisons.
\FloatBarrier

\subsection{Distribution des réplications et stabilité numérique}
La figure~\ref{fig:hist-bootstrap} représente les 2\,000 écarts piétons-conducteurs.
Elle décrit une distribution de statistiques répliquées, et non une distribution de
blessures ou une population de 2\,000 usagers. Les traits encadrent les bornes percentiles.
\figurerapport[0.90]{05_02_distribution_bootstrap_pietons.png}{Distribution bootstrap de l'écart piétons-conducteurs. Chaque valeur résulte d'un nouveau tirage d'accidents dans les quatre années.}{fig:hist-bootstrap}

Les bornes calculées avec les 1\,000 premières réplications sont comparées à celles
des 2\,000 réplications, pour les écarts en points. La variation maximale est de
0,0350 point, soit environ 0,04 point de pourcentage. Ce contrôle indique une faible
sensibilité numérique sur cette exécution~; il ne prouve ni la couverture réelle
des intervalles ni la validité des hypothèses.
\begin{table}[htbp]
\centering\small
\renewcommand{\arraystretch}{1.18}
\caption{Variation maximale d'une borne d'écart entre 1 000 et 2 000 réplications.}\label{tab:stabilite}
\begin{tabularx}{\textwidth}{@{}LLL@{}}
\toprule
Comparaison & Cible principale (points) & Décès seuls (points) \\
\midrule
75+ / 35-44 ans & 0,0350 & 0,0307 \\
Piétons / conducteurs & 0,0118 & 0,0031 \\
Nuit sans éclairage / jour & 0,0161 & 0,0109 \\
Hors / en agglomération & 0,0048 & 0,0045 \\
\bottomrule
\end{tabularx}
\sourcecalcul{notebook 05 ; contrôle numérique sur une graine fixée}
\end{table}

\FloatBarrier

\clearpage
\section{Lecture annuelle et sensibilité aux décès seuls}
\subsection{Les différences gardent-elles le même sens chaque année~?}
Les comparaisons sont recalculées pour chaque année à partir des proportions et des
réplications propres à cette année. Les années elles-mêmes ne sont pas rééchantillonnées.
Les quatre écarts observés restent positifs en 2019, 2020, 2021 et 2022. L'écart d'âge
est cependant plus faible en 2021-2022 qu'en 2019-2020.
\figurerapport{05_03_contrastes_annuels.png}{Écarts annuels pour la cible principale. Les zones représentent des intervalles marginaux par année, sans couverture simultanée ni prévision d'une année future.}{fig:contrastes-annuels}

La répétition du sens des écarts renforce leur constat descriptif sur la période.
Elle ne sépare pas les contributions des autres variables et ne démontre pas un effet
du temps. Les bandes ne sont pas des intervalles de prévision pour 2023 ou 2024.
\FloatBarrier

\subsection{Changer l'événement : décès seuls}
La cible complémentaire vaut 1 pour les décès et 0 pour les trois autres modalités
connues. Les gravités inconnues restent exclues du dénominateur. Les mêmes 467\,079
usagers à gravité connue et les mêmes tirages sont donc utilisés. On compte
11\,996 décès, soit 2,57\,\%.
\begin{table}[htbp]
\centering\small
\renewcommand{\arraystretch}{1.18}
\caption{Écarts et rapports : décès seuls, 2019-2022.}\label{tab:deces_seul}
\begin{tabular*}{\textwidth}{@{\extracolsep{\fill}}lrrrr@{}}
\toprule
Comparaison A / R & Écart (points) & IC à 95\% & Rapport & IC à 95\% \\
\midrule
75+ / 35-44 ans & 6,44 & [6,06 ; 6,84] & 4,27 & [4,00 ; 4,61] \\
Piétons / conducteurs & 2,39 & [2,16 ; 2,63] & 1,98 & [1,88 ; 2,08] \\
Nuit sans éclairage / jour & 4,12 & [3,87 ; 4,37] & 2,87 & [2,74 ; 3,01] \\
Hors / en agglomération & 3,27 & [3,15 ; 3,38] & 3,44 & [3,30 ; 3,58] \\
\bottomrule
\end{tabular*}
\sourcecalcul{notebook 05 ; intervalles percentiles marginaux, 2 000 réplications d'accidents stratifiées par année}
\end{table}

\figurerapport{05_04_contrastes_deces_seuls.png}{Écarts et rapports pour les décès seuls. Une différence absolue plus faible peut coexister avec un rapport plus élevé.}{fig:forest-deces}
\FloatBarrier

Chez les 75 ans ou plus, la proportion de décès est de 8,41\,\%, contre 1,97\,\% chez
les 35-44 ans~: l'écart vaut 6,44 points et le rapport 4,27. Le rapport principal
était de 2,39. Les rapports pour la lumière et l'agglomération sont aussi plus élevés
avec les décès seuls, alors que celui des piétons reste proche de deux.

La définition de l'événement change donc la lecture quantitative. Cette analyse
permet d'examiner la sensibilité à l'hospitalisation~; elle ne justifie pas de choisir
une cible parce qu'elle produit un contraste plus spectaculaire.
\FloatBarrier

\clearpage
\section{Bilan de la première partie et limites}
La première partie a produit une table métropolitaine traçable, des contrôles de
jointure, un portrait des usagers et des comparaisons dont les dénominateurs sont
explicites. Les écarts observés selon l'âge, la catégorie d'usager, la lumière et
l'agglomération persistent dans la lecture annuelle de 2019-2022. Leur amplitude
dépend de la définition de la gravité.

Trois limites restent structurantes. Les accidents enregistrés ne représentent pas
tous les trajets ni nécessairement tous les accidents réels. L'hospitalisation et
les informations inconnues sont sensibles à la collecte. Enfin, les comparaisons
sont brutes et certaines caractéristiques se recoupent. Le bootstrap ne supprime
aucune de ces limites et n'intègre pas toutes les dépendances possibles entre accidents.

La seconde partie étudie plusieurs caractéristiques simultanément et évalue des
probabilités sur une année ultérieure, avec une référence constante et un protocole
temporel déclaré. La gravité et ses dérivés sont exclus des prédicteurs, les traitements
appris restent limités à l'apprentissage et les mécanismes de valeurs manquantes
font l'objet de diagnostics spécifiques.
